% Generated draft: requires mathematical and source review.
\documentclass[11pt]{article}
\usepackage[margin=25mm]{geometry}
\usepackage{fontspec}
\setmainfont{Latin Modern Roman}
\usepackage{amsmath,amssymb,mathtools,unicode-math}
\setmathfont{Latin Modern Math}
\usepackage{enumitem,needspace,xurl,hyperref}
\setlength{\parindent}{0pt}
\setlength{\parskip}{5pt}
\setlength{\emergencystretch}{3em}
\newcommand{\N}{\mathbb{N}}
\newcommand{\Z}{\mathbb{Z}}
\newcommand{\Q}{\mathbb{Q}}
\newcommand{\R}{\mathbb{R}}
\newcommand{\C}{\mathbb{C}}
\newcommand{\F}{\mathbb{F}}
\newcommand{\A}{\mathbb{A}}
\newcommand{\PP}{\mathbb{P}}
\newcommand{\E}{\mathbb{E}}
\newcommand{\eps}{\varepsilon}
\newcommand{\dd}{\,\mathrm{d}}
\begin{document}
\section*{UnsolvedMath ID: 30001719}
\textbf{Source problem code:} catalogue entry
\subsection{Definitions and mathematical statement}
For a specified algebra \(A\), the notation \(\operatorname{proj}A\) conventionally denotes the category of finitely generated projective \(A\)-modules. Its Grothendieck group \(K_0(\operatorname{proj}A)\) is generated by isomorphism classes \([P]\), with relations \([P\oplus P']=[P]+[P']\). The supplied record does not specify \(A\), its base, or the module-side convention.

The retained source extraction concerns a map
\[
I:K_0(\operatorname{proj}A)\longrightarrow\mathcal A_Q^+.
\]
It asserts that this map is canonical and commutes with mutations, and that, in an unspecified example, its image is the generic basis. The extraction further asserts that these results establish the Fock--Goncharov conjecture in that example.

\textbf{Unresolved formulation.} The record supplies neither the quiver \(Q\), the meaning of \(\mathcal A_Q^+\), the construction of \(I\), the relevant mutations and their domains, the generic basis, the example, nor the particular formulation of the Fock--Goncharov conjecture. Consequently, neither a precise mutation-compatibility identity nor the claimed basis identification can be formulated from this evidence. The assertion about the conjecture is retained solely as a supplied source claim, with its scope restricted to the unidentified example.

The supplied statement-verification note says that the official report presents this extraction as a theorem rather than an open problem, and that no separate clean statement is promoted. The registry nevertheless labels the entry \texttt{open} and its statement \texttt{unrecoverable}. These are supplied curation labels, not independently verified mathematical conclusions. No standalone open problem can be recovered from the supplied record.
\subsection{Short English statement}
The record describes a map from a Grothendieck group that reportedly respects mutations and yields a generic basis in a particular example, where a Fock--Goncharov conjecture is reportedly proved. The objects, mutations, example, and conjecture formulation are missing. The supplied verification identifies the passage as a theorem; its registry label \texttt{open} does not specify an open mathematical question.
\subsection{Sources}
\begin{itemize}
\item Supplied local UnsolvedMath registry record, permanent ID \(30001719\), exact source problem code \texttt{catalogue entry}, titled ``Theorem: Mutation Compatibility of Generic-Basis Maps.'' The record belongs to the Oberwolfach Reports Open Problems collection and gives the year 2011; no individual proposer is supplied.
\item \emph{Representation Theory of Quivers and Finite Dimensional Algebras} (2011), Oberwolfach Reports, Report 10/2011, pp.\ 523--608, DOI: \url{https://doi.org/10.4171/owr/2011/10}. This is both the supplied source citation and the supplied literature-source link. No location within the report identifying the particular theorem or example is supplied.
\item Literature assessment dated exactly \texttt{2026-08-21}: supplied dated curation, not independently verified here. It labels the entry open under the stated curation policy because the extraction is an incomplete reference to a theorem and a conjecture, with no identifiable standalone mathematical claim to verify. The background identifies this assessment as literature triage derived from supplied Oberwolfach curation and calls for independent verification of time-sensitive status and cited claims.
\item The supplied statement-verification note says that the official report states the extraction as a theorem rather than an open problem; the original record is retained for provenance, without a promoted clean statement. This description has not been independently checked here.
\end{itemize}
\end{document}
